\section{Historical validation studies and their role in v5}
\label{sec:v5-historical-validation}

This record groups the earlier development studies by the question they answer.
Section~\ref{sec:toys} retains the full-refit validation and support studies;
Section~\ref{sec:selected-observed-results} gives the current observed limits.

\subsection{Kernel, window and normalization development: June--August 2026}
Early window and cross-kernel studies compared local polynomial backgrounds
with GP interpolation. Later work widened the fit support beyond the search
interval, adopted the moving $\pm2.25\sigma_m$ exclusion, and retained
full-template signal normalization in the local likelihood. These studies
explain the current geometry; their old limit illustrations are historical.

The v4.2 follow-up separated display repairs from conditional window
replacement and shared-coupling tests. Replacing selected counts creates
neither a new background-only observation nor independent signal evidence.
Appendix~\ref{sec:candidate-removal} extends that question with declared
whole-bin interventions and comparison fits.

\subsection{Source, exposure and threshold tests: August 2026}
The v4.5--v4.8 program first used small pilot ensembles, then 100-background
full-refit ensembles, to separate the fitted source shape from the imposed
exposure. Its matched-reference injections use each background toy's reference
uncertainty, fixed before adding the signal. The displayed results in Section~\ref{sec:toys}
retain the original ensemble definitions and their conditional nature.

Threshold-focused tests showed why passing a smooth-source ensemble does not
establish performance near every steep turn-on. The v4.9 qualification and
v4.9.5 support continuation examined several lower edges. All eligible edges
failed the first, more stringent screen. A later practical screen was fixed
before the held-out continuation and selected 36~MeV within the retained tie
rule. The surviving under-recovery at 55~MeV is part of the result, not an
exception removed from the figure. That sequence explains both the adopted
support and the limits of its validation.

\subsection{The 2016 support and fit-reproducibility checks: September 2026}
The v4.9.7 study tested lower support edges 28--34~MeV at 44, 49, 54 and
59~MeV using 25 paired backgrounds and matched 0-, 2- and 5-reference-sigma
injections. Every edge violated the predeclared gross-bias guard. The rule
returned no provisional edge, so the proposed 75-background continuation and
new combination did not follow from that study. Its broad-tail generating
component carried a failed source-fit flag and was retained only as an
alternative background for the stress test.

Later v4.9.9--v4.9.11 studies checked whether a calculation run in a
separate process reproduced the fitted GP parameters for the 30--210~MeV
support. Agreement within the required $10^{-6}$ tolerance was obtained at
87 of 142 masses; 55 did not meet that tolerance. Recalculating the
background predictions from the saved parameters, without rerunning the
optimizer, passed the subsequent numerical checks. The v4.9.12 results
therefore use those saved parameters. The discrepancies in the independent
parameter fits remain unresolved and are separate from the failed
background-validation tests described above.

Full-refit closure, frozen-GP expected limits, empirical \CLs{} calibration
and coherent global scans test different parts of the procedure. More draws
in one cannot replace missing validation of another. This distinction allows
reproducible observed curves to coexist with open physics-review questions.
